% !TeX root = ../../paper.tex

\subsection{Scenario and Game Flow}
\label{sec:apparatus-scenario}

The player starts inside a dimly lit wooden cabin. A lock with four coin-shaped slots bars the front door; the way out is to find four golden coins hidden around the room and insert them into the lock. Dark ambient music and the cramped, cluttered interior set the tone before anything threatens the player.

Each coin hides inside one of four challenges, described in Section~\ref{sec:apparatus-challenges}, so the search forces the player to interact with every station in the room. The game has no time limit and no fail state; the discomfort lies in what the player has to touch to get each coin. Once the fourth coin clicks into place, the lock opens and a win screen ends the session. We chose this flow for the study because it guarantees that each participant encounters all four haptic (or non-haptic) interactions in a single playthrough, while leaving the order of the challenges up to the player.

A game manager script drives this loop. The same scene runs in both study conditions; the conditions differ in whether the physical proxies lie in the tracking space at the positions of their virtual counterparts.
